\section{Gauge and the Photon}

The force field of the vibe mesh is \textbf{one tone}, not a second field. Everything in this section follows from a single commitment that the rest of the paper has carried from the start. There is one substance, the vibe, and the only value it takes in a site is one tone in $\{-1, 0, +1\}$, read as pain, peace, and pleasure. Matter is that tone bound. The photon is that same tone waving. Nothing else is added.

This is a strong claim, and it is worth saying plainly before any machinery appears. Electromagnetism, in the usual account, is a field of its own, distinct from the matter it pushes around. Here it is not. The light that crosses a room and the electron that sits in an atom are two motions of the very same tone field. The aim of this section is to show that this single field is enough. It gives a local $U(1)$ symmetry, a discrete Gauss law that was measured, a Coulomb $1/r$ potential, the Lorentz force, and a non-abelian extension reaching all the way to a grand-unified group. Along the way it records one wrong turn, a second ternary field that was worked out in full and then abandoned, kept only so it is never tried again.

\begin{principle}[One tone, two motions]
There is one ternary tone. Matter is the tone bound. The photon is the tone waving. That single fact is mass-energy equivalence stated in the substrate.
\end{principle}

\subsection{The one-tone picture}

Recall the local structure of a dock. A \emph{dock} is one cell of the $\{3,4,3,4\}$ mesh, a here-now through which tones pass. Out of each dock run exactly $24$ \emph{sites}, the $24$ directions, and these sites are the $D_4$ root system. There is no rest slot, no twenty-fifth direction, no home. A tone lives in a site and streams along it each beat.

The symmetry of the $24$ sites splits them into three eights. This is \emph{triality}, the order-three outer symmetry of $\mathrm{SO}(8)$ that permutes the three pieces. The decomposition is $8v + 8s + 8c$, a vector plus two spinors, and it is the spine of everything that follows.

\begin{table}[ht]
\centering
\begin{tabular}{@{}lll@{}}
\toprule
sector & size & role \\
\midrule
$8v$ (vector) & $8$ & the photon \\
$8s$ (spinor) & $8$ & one fermion chirality \\
$8c$ (spinor) & $8$ & the other fermion chirality \\
\bottomrule
\end{tabular}
\caption{Triality splits the $24$ sites of a dock into three eights. The vector eight is the photon. The two spinor eights are the two chiralities of matter.}
\label{tab:triality-sectors}
\end{table}

The two spinor eights $8s$ and $8c$ are the fermion fields, treated in the section on spin and fermions. The vector eight $8v$ is the photon. They are not three fields. They are three faces of one tone field, related by triality.

What separates matter from the photon is not the substance but the motion. Both are excitations of the one tone field. The difference is whether the excitation is bound or free, and whether it carries net charge.

\begin{table}[ht]
\centering
\begin{tabular}{@{}llll@{}}
\toprule
 & what it is & charge & motion \\
\midrule
matter & a localized, bound charge & net non-zero & trapped, can be at rest \\
photon & a free, propagating wave (a phonon) & net zero, $+$ and $-$ balance & always moving, massless \\
\bottomrule
\end{tabular}
\caption{Matter and the photon are two motions of the same tone field. Matter is bound and charged. The photon is free and neutral.}
\label{tab:matter-vs-photon}
\end{table}

This is $E = mc^2$ read off the substrate. Rest mass is trapped energy. A bound circulating excitation has mass because its energy is confined to a region and cannot leave. Un-trapping that excitation turns the mass into free radiation. One substance, two states, bound and free.

\begin{definition}[The charge]
The \emph{charge} of a dock is the signed sum of its tones, an integer fixed by the \emph{lean}, the value order $-1 < 0 < +1$. Matter is a region of net non-zero charge. A photon is a net-zero ripple, its plus and minus tones balancing exactly.
\end{definition}

\subsection{The drumhead}

The cleanest way to feel why one field suffices is the drumhead. Strike a drumhead once. It vibrates, radiates sound into the surrounding air, and then goes still.

The air is not a second field. It is the same medium the drumhead sits in, and sound is just that medium's waves. The disturbance leaves the drumhead, travels out as sound, and never comes back, so the drumhead settles to rest.

Here the medium is the one tone field.

\begin{itemize}
\item Matter is a localized excitation of it, the drumhead.
\item The photon is a wave of it, the sound.
\item The \emph{bath} is the open mesh. The mesh grows at its \emph{wake}, its ever-expanding frontier, so radiation falls behind the wake and never returns.
\end{itemize}

The wake is doing real work here. Because the mesh is open and forever growing at its edge, an outgoing wave has somewhere to go and nothing to bounce off. That is what lets a disturbed body shed its excess and come to rest. A closed box would reflect the sound back. An open, growing mesh does not.

\subsection{Why one field suffices}

Everything follows from the one tone and the one charge. Three consequences make the single-field picture self-consistent.

\begin{itemize}
\item \textbf{One charge.} The only conserved quantity is the $U(1)$ charge, the signed sum of tones in a dock. Matter is a region of net charge. A phonon is a net-zero ripple. There is no second charge, so there is no second field for it to source.
\item \textbf{Neutral radiation.} A phonon carries equal plus and minus, so it ferries the disturbance away without carrying off any net charge. The body sheds the disturbance while keeping its identity, its charge, intact.
\item \textbf{Settling.} Perturb the body's motion with a charge-neutral kick. It radiates that momentum away as a phonon, the bath absorbs it at the wake, and the body returns to its ground state.
\end{itemize}

That last point is the corrective agency of a \emph{self}, an emergent bound pattern, expressed in one field. A self can absorb a hit and recover precisely because it can hand the excess to the radiation channel, and that channel is not a separate ingredient. It is a wave of the same tone the self is made of.

\begin{result}[The phonon is emergent, not per-cell]
A soft, gapless sound mode does emerge from the discrete ternary tones on the $\{3,4,3,4\}$ substrate, with frequency growing linearly with wavevector at constant speed, $\omega = c\,k$, built and measured with no real numbers \cite{pollard2026vibetest}. The ternary tone is too coarse to carry a soft ripple per cell, so softness is a long-wavelength collective property of many tones, not a small-amplitude property of one. The photon is a coarse-grained excitation of the one tone field.
\end{result}

This deserves a careful word, because it corrects a tempting mistake. The tone has only three values, so the smallest single-cell change is order one. There is no gentle per-cell ripple. For a while this looked fatal, as if no soft radiation could exist on a ternary substrate at all. The resolution is that softness is long \emph{wavelength}, not small \emph{amplitude}. A long, slow wave spread across many tones is soft even though each tone flips by a whole unit. Boolean and ternary lattice gases recover ordinary sound exactly this way. So the phonon is real, it just lives one layer up, as a collective mode, never as a base per-cell object.

\subsection{Local $U(1)$ and the Gauss law}

The \emph{knit}, the collide-then-stream law, conserves charge \textbf{locally}. Charge in equals charge out, per dock, every beat. The tones in a dock collide by a fixed reversible table that preserves their signed sum, then stream one dock along their sites. Nothing creates or destroys net charge anywhere.

A local conservation law is exactly a Gauss law. This is not an analogy, it is the same statement. A global conservation law says the total is fixed. A local one says it is fixed dock by dock, which is the discrete divergence condition that the conserved current sources a field.

\begin{principle}[Local conservation is a Gauss law]
Charge conserved per dock per beat is a discrete Gauss law. The conserved current sources an emergent $U(1)$ field. This holds exactly, on every substrate, because it is a property of the knit and not of the geometry.
\end{principle}

\subsection{The emergent gauge field}

To check that the emergent field is a genuine gauge field, put a vortex gauge field on the mesh. Each site link carries a phase, and the phases wind around a plaquette to enclose a flux $\Phi$. A true gauge theory has physical content only in the gauge-invariant field strength, never in the link phases themselves. The measurements confirm exactly this.

\begin{table}[ht]
\centering
\begin{tabular}{@{}ll@{}}
\toprule
observable & result \\
\midrule
Wilson loop, sum of link phases around a loop & equals the enclosed flux $\Phi$ \\
gauge invariance & unchanged under $A \to A + d\lambda$ \\
Aharonov-Bohm phase, charge encircling the flux & exactly $\Phi$, independent of loop and gauge \\
non-enclosing loop & zero \\
\bottomrule
\end{tabular}
\caption{The emergent field is a genuine abelian gauge theory. Its physical content is the gauge-invariant flux, not the link phases.}
\label{tab:gauge-observables}
\end{table}

The Wilson loop returns the enclosed flux. A loop that encloses no flux returns zero. A charge carried around the flux picks up exactly the Aharonov-Bohm phase $\Phi$, and that phase does not depend on the path or on the choice of gauge. These are the defining marks of an abelian gauge theory. The emergent $U(1)$ is electromagnetism.

\begin{result}[A genuine $U(1)$ gauge theory]
The emergent field built from the conserved tone current is gauge-invariant, its Wilson loop equals the enclosed flux, and the Aharonov-Bohm phase is exactly $\Phi$ independent of loop and gauge \cite{pollard2026vibetest}. The substrate carries electromagnetism as an emergent $U(1)$ gauge theory.
\end{result}

\subsection{The Coulomb potential}

The static sector of this field is the flat three-dimensional Laplacian Green's function. That function is the Coulomb $1/r$ potential. It was obtained by the difference method and is set out in full in the section on electromagnetism.

\begin{result}[The Coulomb $1/r$]
The static emergent field is the flat-$3$D Laplacian Green's function, the Coulomb $1/r$ potential, recovered by the difference method \cite{pollard2026vibetest}.
\end{result}

\subsection{The Lorentz force and lattice QED}

The matter sector and the photon sector are not two worlds evolving side by side. Under one coupled knit they act on each other, and they do so both ways. A moving charge in the fermion sector \textbf{sources} a radiating field in the photon sector, whose energy grows from exactly zero. A background field in the photon sector \textbf{accelerates} the fermion at the Newton rate.

That deflection of a charge by a field is the Lorentz force, and it was measured under one coupled knit. The discriminating control is the single coupling constant. At zero coupling the one knit decouples and both effects vanish exactly, no field is sourced and no acceleration occurs. So the sectors are bound by the knit, and the binding switches off together. The full dynamical account is in the section on photon-fermion co-emergence.

\begin{result}[The sectors couple dynamically, both ways]
Under one coupled knit a moving charge sources a radiating field whose energy grows from zero, and a background field accelerates the fermion at the Newton rate \cite{pollard2026vibetest}. At zero coupling both effects vanish exactly. The deflection is the Lorentz force.
\end{result}

The whole construction is lattice quantum electrodynamics in the discrete sense. The matter sites carry the charge. The site links carry the gauge field. The Wilson loop and the Gauss law are the lattice observables. There is no continuum input anywhere, only the knit on the mesh.

\subsection{The non-abelian extension}

The same Wilson-loop construction, run with matrices in place of phases, gives a non-abelian gauge field. It was verified with $\mathrm{SO}(3)$ link matrices, a prototype whose structure is identical to the grand-unified case.

\begin{table}[ht]
\centering
\begin{tabular}{@{}ll@{}}
\toprule
property & result \\
\midrule
Wilson loop, trace of the ordered product & gauge-invariant \\
holonomy & non-trivial, a genuine curvature \\
link matrices & do not commute, non-abelian \\
\bottomrule
\end{tabular}
\caption{Replacing link phases by link matrices yields a non-abelian gauge field, verified with $\mathrm{SO}(3)$ matrices.}
\label{tab:nonabelian}
\end{table}

A non-abelian gauge group reduces to the symmetry of the collide table. Gauging the larger group is nothing more than the existence of a collision invariant under it, which the substrate supplies through the $D_4$ and $F_4$ symmetry of the $24$ sites. The grand-unified group is taken up in the section on the standard model.

\begin{principle}[Gauging is a collision symmetry]
A non-abelian gauge group reduces to the symmetry group of the collide table. Gauging it is the existence of a collision invariant under that group, which the substrate's site symmetry provides.
\end{principle}

\subsection{The wrong turn, a second field abandoned}

Honesty requires recording a route that was followed in detail and then dropped. An earlier line of reasoning concluded that the substrate needed a \textbf{second ternary field}, a $\mathbb{Z}_3$ gauge field, to carry the photon. That conclusion is wrong. It violates the founding commitment that there is one tone, not two.

The argument ran in three steps.

\begin{enumerate}
\item Emission needs a neutral photon. This is correct.
\item No fixed sector of the $24$ sites can be a static neutral photon, by the no-coloring enumeration. This is also correct.
\item Therefore the photon must be a separate neutral field. This is wrong.
\end{enumerate}

The flaw is the third step. The fact that no fixed sector of sites is the photon rules out the photon being a \emph{static subset} of directions. It says nothing about the photon being a \emph{dynamical wave} of the one tone.

\begin{principle}[A wave is not a sector]
The no-coloring result forbids the photon from being a static subset of the $24$ sites. A wave is not a subset of sites, so the result does not touch it. The photon is a dynamical excitation of the one tone, and the second field is unnecessary.
\end{principle}

The full $\mathbb{Z}_3$ specification was worked out completely, two fields per dock, an exact emission involution table, and a discrete Gauss law of its own. It is consistent. It is simply richer than needed, a heavier design where the one-tone wave already does the work. It is kept only for completeness in the appendix on abandoned routes.

\subsection{What this keeps}

The one-tone account keeps the base exactly as it was. Nothing is added to the eight atoms.

\begin{itemize}
\item One tone, ternary in $\{-1, 0, +1\}$.
\item The $24$ sites of a dock, the $D_4$ directions, with no rest slot.
\item The reversible knit, collide then stream.
\item The beat, the wake, the open mesh, and the emergent arrow.
\end{itemize}

No second field. No second alphabet. No real numbers at the base. Matter and the photon are not new ingredients. They are two excitations of the one tone field, bound and free.

The whole $U(1)$ gauge sector comes from the one tone. Local conservation forces the Gauss law. The wave is the photon. The static field is the Coulomb $1/r$. The coupled knit gives the Lorentz force. The matrices give the non-abelian extension. The second-field route was a wrong turn, recorded plainly so it is not tried again.
